\documentclass[11pt,twoside,openright]{book}

\usepackage[T1]{fontenc}
\usepackage[utf8]{inputenc}
\usepackage{amsmath,amssymb,amsthm,mathtools,bm}
\usepackage{mathpazo}
\usepackage{microtype}
\usepackage[inner=1.35in,outer=1.1in,top=1.25in,bottom=1.35in]{geometry}
\usepackage{booktabs}
\usepackage{longtable,array}
\usepackage{fancyhdr}
\usepackage[hidelinks]{hyperref}
\usepackage{cleveref}

\setlength{\parindent}{1.4em}
\setlength{\parskip}{0pt}
\linespread{1.07}
\emergencystretch=3em

\pagestyle{fancy}
\fancyhf{}
\fancyhead[LE]{\small\thepage\quad\textsc{A Cinema of Individual Editions}}
\fancyhead[RO]{\small\nouppercase{\leftmark}\quad\thepage}
\renewcommand{\headrulewidth}{0.3pt}
\fancypagestyle{plain}{\fancyhf{}\fancyfoot[C]{\small\thepage}\renewcommand{\headrulewidth}{0pt}}

\theoremstyle{plain}
\newtheorem{theorem}{Theorem}[chapter]
\newtheorem{proposition}[theorem]{Proposition}
\newtheorem{lemma}[theorem]{Lemma}
\newtheorem{corollary}[theorem]{Corollary}
\theoremstyle{definition}
\newtheorem{definition}[theorem]{Definition}
\newtheorem{example}[theorem]{Example}
\theoremstyle{remark}
\newtheorem{remark}[theorem]{Remark}

\newcommand{\Capac}{\mathrm{Cap}}
\newcommand{\Rob}{\mathrm{Rob}}
\newcommand{\Inv}{\mathrm{Inv}}
\DeclareMathOperator{\loss}{loss}

\title{\Huge\bfseries A Cinema of Individual Editions\\[1.2ex]
\large Appendix C: Notation and Relation to the Earlier Volume\\[0.4ex]
\large Working Draft}
\author{Flyxion\\[0.3ex]\small Independent Researcher}
\date{November 2026}

\begin{document}

\frontmatter
\maketitle
\tableofcontents
\mainmatter
\appendix
\setcounter{chapter}{2}

\chapter{Notation and Relation to the Earlier Volume}
\label{app:notation}

This appendix does three things. It tabulates the notation of the book, including the places where one symbol carries two meanings. It sets out the correspondence between the objects of the earlier volume, \emph{Generative Film: Categories, Sheaves, and Types for Cinema}~\cite{flyxion2026gf}, and the objects of this one. And it lists the points at which the two books use different words for the same thing, or the same word for different things. Where the correspondence is a reading and not a theorem, the appendix says so.

\section{Notation of This Book}\label{sec:notation}

The table gives the symbols in order of first use. The chapter column gives where the symbol is introduced or most used.

\begin{longtable}{@{}>{\raggedright\arraybackslash}p{2.6cm}>{\raggedright\arraybackslash}p{8.4cm}>{\raggedright\arraybackslash}p{1.6cm}@{}}
\toprule
Symbol & Meaning & Chapter\\
\midrule
\endfirsthead
\toprule
Symbol & Meaning & Chapter\\
\midrule
\endhead
\bottomrule
\endfoot
$h^\ast$ & the script history of a scene & 1, 3\\
$m$ & number of slots in a scene plan & 1\\
$t_j$, $V_j$ & story time and set of strips in frame at slot $j$ & 1\\
$r_j$ & required appearance at slot $j$, mown or unmown for each strip in frame & 1\\
$F(s)$ & footprint of strip $s$, the set of slots that show it & 1\\
$A$, $B$, $C$ & the three productions of the lawn scene, with audiences $60$, $30$ and $2$ & 1\\
$D$ & the direction rule on adjacent slots & 1, 4\\
$G$, $\alpha_a,\alpha_b,\alpha_c$ & the agreement rule, and its restrictions to the strips $a$, $b$, $c$ & 8, 9\\
$\omega$ & the take-level amendment withdrawing a take variant & 9\\
$K$ & number of distinguishable editions; in Chapter~14 a viewer's knowledge & 1, 14\\
$B$ & storage cost of footage, rules and metadata; also the second production in tables & 1\\
$\rho$ & the rate $\log_2K/B$; in Chapter~15 a neighborhood radius & 1, 15\\
$\varepsilon$ & tolerance of perceptual equivalence & 2\\
$M_\varepsilon$, $N_\varepsilon$ & packing and covering numbers at tolerance $\varepsilon$ & 2, A\\
$\pi$, $\Capac$ & projection of realizations to class tuples; capacity $\lvert\pi(X)\rvert$ & 9\\
$\Rob$, $\Inv$ & robust and invariant capacity of a family of amendments & 9\\
$U$, $U_p$ & size of an audience; of the audience of locale $p$ & 1, 11\\
$r^\ast$ & least radius of a repair & 10\\
$G_p$ & set of global classes that locale $p$ can realize & 11, A\\
$\ell$, $s$ & number of slots with a local trace, and of shared slots, in Chapter~11 & 11\\
$N$, $w$, $h$, $y$ & viewings, hours per week, hours per viewing, years, in Chapter~13 & 13\\
$H$, $C_f$ & set of histories of the engine scene; the set a film $f$ commits to & 3, 14\\
$R(K)$ & the readings of the histories in a knowledge state $K$ & 14\\
$X$ & a reference corpus in Chapter~15; a set of realizations in Chapter~9 & 9, 15\\
$M_j$, $D_j$ & compatibility matrix and weight matrix of slot $j$ & A\\
$s_j(k)$ & number of admissible completions from slot $j$ beginning with class $k$ & A\\
\end{longtable}

Five of these symbols are reused. The letter $\rho$ is the rate in Chapter~1 and the neighborhood radius in Chapter~15, and the chapters do not meet. The letter $K$ counts editions in Chapter~1 and names a knowledge state in Chapter~14. The letter $N$ is the number of viewings in Chapter~13, and the covering number $N_\varepsilon$ carries a subscript. The letter $G$ is the agreement rule in Chapters~7 to~10 and the set $G_p$ of global classes of a locale in Chapter~11, the latter always subscripted. The letter $M$ is the packing number when subscripted by a tolerance and the compatibility matrix of a slot in Appendix~A when subscripted by a slot. The letter $B$ is a storage cost in Chapter~1 and a production only inside the tables of the lawn scene. None of these meanings is used in the same passage as another.

\section{Correspondence of Objects}\label{sec:corr}

The earlier volume treats film as a composition problem and develops category, sheaf, and type theory for it. This book is about a different question, how many variants of a film a production can sustain, and it uses little of the earlier machinery directly. The correspondences are of three kinds: an object of this book is an instance of an object of the earlier volume, a result of this book can be read as a result of the earlier one, or the two stand in a looser relation. The table says which.

\begin{longtable}{@{}>{\raggedright\arraybackslash}p{3.7cm}>{\raggedright\arraybackslash}p{3.7cm}>{\raggedright\arraybackslash}p{5.1cm}@{}}
\toprule
Earlier volume & This book & Relation\\
\midrule
\endfirsthead
\toprule
Earlier volume & This book & Relation\\
\midrule
\endhead
\bottomrule
\endfoot
Layered world: history $A\subseteq E^\ast$, fold $M:A\to Q$, appearance $R:Q\to I$ (Chapter 11) & Script history $h^\ast$, set of mown strips $M(h^\ast_{\le t})$, picture of the strips in frame at a slot & Instance. The lawn scene is a layered world whose state is the set of mown strips.\\
Authority principle (Principle 11.5): history authoritative at commitment, appearance never authoritative & Ledger and archive of Chapter 12: records of assignments and takes are kept, rendered frames are not & Reading. The ledger stores the tuple from which the appearance is recomputed.\\
Future equivalence $\sim$ and the quotient $A/{\sim}$, the minimal state space (Proposition 11.2) & Future-equivalence quotient used to count overcapture by noise in Chapter 3 & Instance.\\
Several folds of one history: physical, semantic, provenance (Section 11.3) & Commitment set of an edition, readings of Chapter 3 & Reading. A commitment set is a semantic fold of the histories that agree with what an edition shows.\\
Admissible histories, prefix-closed $A$ (Chapter 9, 11, 12) & Admissible class tuples under the direction and agreement rules & Analogy, one level up. Admissibility here is of tuples of slot classes, not of event words.\\
Viewer classes $V$, schedule $\mathbb N=\bigsqcup N_v$, rate budget $\lvert V\rvert\le\lfloor r/r_{\min}\rfloor$ (Proposition 13.1) & Locales and audiences, Hall condition (Chapter 11, Appendix A) & Different resource. The earlier budget shares frames of one screen, and this book counts editions and says nothing about how they are delivered.\\
Joint constraint $J\subseteq\prod_vF_v$, subsheaf if and only if a product (Proposition 13.2) & Worldwide distinctness of editions, a non-product constraint on the editions of all viewers & Reading. Each viewer's edition can be admissible and the tuple of editions not, which is the failure of Hall's condition in Chapter 1.\\
Branching film, tree $T$ with $L$ leaves, shared segments (Section 13.3, Principle 13.4) & Shared slots among locales (Chapter 11) & Reading. A shared slot cannot carry a local trace, as a shared segment cannot require a witness from one branch.\\
Version, situation $\theta$, dissimilarity $d_\theta$, three response variables (Section 14.1) & Edition, the detection, comprehension and affect dissimilarities (Chapter 2) & Instance. The situation $\theta$ is left unspecified in this book.\\
Thresholded similarity $\approx_\varepsilon$ is not transitive (Proposition 14.3), drift $\le n\varepsilon$ & Non-transitivity of $\varepsilon$-similarity, packing and covering numbers (Chapter 2) & Instance, with a counting form added.\\
Effective dimension $d_{\mathrm{eff}}$ from singular values of a Jacobian (Proposition 14.6) & Packing number $M_\varepsilon$ and covering number $N_\varepsilon$ & Counterpart. $d_{\mathrm{eff}}$ counts continuous directions, and the packing number counts distinguishable editions of a discrete pool.\\
Constraints, theory and model operators $\mathrm{Th}$, $\mathrm{Mod}$; adding constraints shrinks $\mathrm{Mod}$ (Chapter 5, Exercise 5.2) & Amendments as subsets $X'\subseteq X$ of realizations (Chapters 8 and 9) & Instance. A correction is a constraint added to the script.\\
Compression as adjunction $D\dashv G$ (Section 5.3) & The rate $\rho=\log_2K/B$ (Chapter 1) & Loose. The rate is a number and not an adjunction. $B$ plays the role of the description.\\
Command language $\mathcal K$, realizations $R(c)$, refinement (Definition 16.3) & Nested commands narrowing slot domains (Chapter 16) & Instance for the narrowing. Prefix-free decoding is not used.\\
Acoustic occupation, contracts, the measured passage (Sections 15.3 and 17.9) & Chapter 6 and Study Three of Chapter 17, with the same data & Same data. This book adds timing editions and a correction to the reference.\\
Register of Claims (Appendix D) & Statements of assumptions beside each result, and Chapter 18 & Different organization. This book has no register, and collects its open questions instead.\\
\end{longtable}

The relations should be read with care. Only the rows marked as instances, which are the layered world, the future-equivalence quotient, the perceptual dissimilarities, the amendments, the commands and the measured passage, claim that an object of this book \emph{is} an object of the earlier one. The rows marked as readings are interpretations that this book offers and does not prove. The row on viewer classes is marked as a different resource because the two books count different things: the earlier one a rate of frames, and this one the number of distinct editions.

\section{Different Words for the Same Thing}\label{sec:words}

The working title of this volume was \emph{A Cinema of Individual Editions}, and the standalone chapter files retain it in their headers; the merged book carries the title \emph{Hypernormal Cinema}.

The words below differ between the two books while the objects are the same or nearly so.

\begin{longtable}{@{}>{\raggedright\arraybackslash}p{3.3cm}>{\raggedright\arraybackslash}p{3.3cm}>{\raggedright\arraybackslash}p{5.9cm}@{}}
\toprule
Earlier volume & This book & Note\\
\midrule
\endfirsthead
\toprule
Earlier volume & This book & Note\\
\midrule
\endhead
\bottomrule
\endfoot
Version & Edition & The earlier word attaches presentation conditions to a film, and this book uses the same idea for the film that a viewer is shown.\\
Viewer class & Audience, locale & The earlier word names a selector worn by viewers, and this book groups viewers by production or locale.\\
Appearance & Take shown at a slot, class at a slot & The earlier word is a function of state, and this book's classes are the equivalence classes of takes under exact slotwise equivalence.\\
Scenario space, film as a path & Space of possible films, edition as a path through the graph of editions & Chapter 4 draws the graph of editions of a scene, a path being a choice of class at each slot.\\
Effective dimension & Packing number, capacity & See the table of correspondences.\\
Tolerance $\varepsilon$ and $\approx_\varepsilon$ & Tolerance $\varepsilon$ & Same symbol, with the books agreeing that thresholded similarity is not an equivalence.\\
Schedule & Cyclic schedule & The earlier word is a partition of frame indices, and this book's is an assignment of ranks to viewers over rounds in Chapter 13.\\
\end{longtable}

Three words are used differently. The word \emph{admissible} describes event words in the earlier volume and tuples of classes here. The word \emph{rate} is a number of frames per second there and the ratio $\log_2K/B$ here. The word \emph{principle} labels each modeling hypothesis in the earlier volume and is not used in this one, where hypotheses are stated in the passage where they enter and collected in Chapter~18. The word \emph{fold} has the same meaning in both, a function from histories to what a history leaves behind, though this book uses it for a commitment set and not for a state.

\section{What the Appendix Establishes and What It Assumes}\label{sec:limits-c}

The tables record the author's account of the two books and have not been checked against anything but the books. The notation table covers the symbols that a reader needs to follow the chapters and does not claim to be complete. The correspondence table is selective: the earlier volume's categories, sheaves, and dependent and linear types are not used in this book beyond the readings in the table, and a reader who wants the connection to those theories will find it in the earlier volume and not here.

\begin{thebibliography}{99}

\bibitem{flyxion2026gf}
Flyxion.
\newblock \emph{Generative Film: Categories, Sheaves, and Types for Cinema}.
\newblock Unpublished monograph, October 2026.

\end{thebibliography}

\end{document}
